\section{Discussion and Limitations}
% \begin{itemize}
%     \item In practice, we may not know if the examiner is accurate or not without privileged info about the environment. 
%     \item Connection with LLM as judge. Inception score. Generative models more broadly. 
%     \item We could run different examiners.
%     \item True environment as an optimal examiner. 
%     \item L
% \end{itemize}
\textbf{Limitations.} The perceived regret is produced by an examiner, which is also a learning agent (that cannot influence the future). When deployed alongside the agent, there are no guarantees that its perceived regret is actually meaningful and it could be wildly incorrect. 
This may be an unavoidable price of seeking to evaluate agents in arbitrary worlds---the general applicability impacts the reliability of the assessment. This issue is akin to the problem of generalization. When an agent encounters a truly novel situation, there cannot be any guarantees and we can only hope it generalizes correctly. 



\paragraph{Models evaluating other models.} The definition of perceived regret is similar in spirit to previously-used performance measures in the machine learning community. Generative models are one area where it has been common to assess the performance of one model with others. 
One example is the Inception score used as a performance metric to measure diversity and fidelity of a model's generated images by passing them through the Inception network. Another is CLIPScore in a similar context for test-to-image generation, utilizing the CLIP model as an evaluator. 
More recently, LLM-as-judge is a common framework used to automatically assess the outputs of a given LLM by feeding it to  multiple other LLMs.
LLMs have also been used to inspect their own outputs, critiquing and refining them to produce higher-quality answers \wes{add cite}. 

\paragraph{The world as an examiner.} When defining the true value and true regret, we 
If we view the agent as part of the world, meaning the world has access to the agent's policy (and learning algorithm), we can view the world as an examiner that outputs $\Vlow_t =V^\pi_t$ and $\Vhi =V^*_t$. 

% \paragraph{Tandem Learning}